Começamos por fixar a notação e os objetos que toda a replicação utilizará. Seja um modelo
de linguagem \textit{transformer} composto por $L$ blocos, de forma que, para um
\textit{prompt} de entrada $X$, o bloco $l$ possui ativações
$A_l(X) \in \mathbb{R}^{t \times k }$, onde $t$ é o comprimento da sequência e $k$ é a
dimensão oculta. Assim como em Beaglehole et al., a característica associada a um
\textit{prompt} é a linha do último \textit{token} real de cada saída de bloco,
$a_l \in \mathbb{R}^k$. Para cada conceito reunimos $n$ \textit{prompts} rotulados
$\{X^{(i)}, y^{(i)}\}$, onde $y \in \{0,1\}$ registra a presença do conceito em cada
\textit{prompt}, passamos cada \textit{prompt} pelo modelo, e empilhamos as ativações de
último \textit{token} por bloco na matriz $Z_\ell \in \mathbb{R}^{n \times k}$ cuja
$i$-ésima \emph{linha} é $\left(a_\ell^{(i)}\right)^{\top}$. Note que extrair um vetor de
conceito não requer texto gerado algum, apenas pares
$(\textit{prompt}, \text{rótulo})$ e uma passagem direta para cada um, o que é o que torna
a coleta de dados barata e o que os estágios seguintes de \textit{steering} e \textit{monitoring}
tomam como base. É sobre esse baixo custo que Beaglehole et al. se apoiam quando extraem um
conceito e lhe aplicam \textit{steering} em menos de um minuto com menos de $500$ amostras.

\subsection{\textit{Recursive Feature Machine}}

A \textit{Recursive Feature Machine} (RFM) de \cite{radhakrishnan2024mechanism} pode ser
lida, informalmente, como uma contraparte supervisionada e não linear da
\textit{Principal Component Analysis} (Análise de Componentes Principais, em tradução
livre, doravante PCA), no sentido
de que, onde a PCA retorna as direções de maior \emph{variância} nos dados, a RFM retorna
as direções às quais um \emph{preditor} ajustado é mais sensível, por meio do
\textit{Average Gradient Outer Product} (Produto Externo Médio dos Gradientes, em tradução
livre, doravante AGOP).
Não se trata de uma generalização da PCA em sentido estrito, uma vez que a PCA não é um
caso particular da iteração abaixo, mas a analogia é a correta para se compreender o que a
matriz $M$ está fazendo.

O preditor, para este caso, é uma \textit{kernel ridge regression} (doravante KRR), isto é,
uma regressão \textit{ridge} com \textit{kernel}, dotada de um \textit{kernel}
Mahalanobis-Laplace parametrizado por uma matriz semidefinida positiva (PSD) $M$ com
largura de banda $\sigma$, dado por

\[
  K_M(x,z) \;=\; \exp\!\left(-\frac{1}{\sigma}\sqrt{(x-z)^{\top} M (x-z)}\right)
\]

onde $\sigma>0$ é a largura de banda do \textit{kernel}. Note que escrevemos $\sigma$ onde
Beaglehole et al. escrevem $L$, de modo a evitar um conflito com o número de blocos do
\textit{transformer} $L$.

Inicializamos $M_0 = I$ e, para $t = 0, \dots, T-1$:

\begin{enumerate}
  \item \textbf{\textit{Kernel ridge regression}.} Ajusta-se
  \[
    \alpha_t = \left(K_{M_t}(Z,Z) + \lambda I\right)^{-1} y \;\in\; \mathbb{R}^{n},
  \]
  \[
    f_t(z) = \alpha_t^{\top} K_{M_t}(Z,z) = \sum_{i=1}^{n} \alpha_{t,i}\, K_{M_t}\!\left(a^{(i)}, z\right).
  \]

  \item \textbf{AGOP.}\footnote{Centramos os gradientes na média antes de formar o produto
externo, conforme pré-registrado em \texttt{plans/GUPPYLM\_GUIDE.md} §1a e conforme
implementado em \texttt{rfm\_guppy/rfm.py}. Sem a centragem, o autovetor principal de
$M_{t+1}$ é puxado em direção ao gradiente médio, e não à direção de maior variação dos
gradientes.} Recalcula-se $M$ como o produto externo médio dos gradientes do preditor em
  relação à entrada, tomados sobre os pontos de treinamento:
  \[
    M_{t+1} = \frac{1}{n}\sum_{i=1}^{n}
      \left(g^{(i)} - \bar{g}\right)\left(g^{(i)} - \bar{g}\right)^{\top},
  \]
  \[
    g^{(i)} = \nabla_z f_t\!\left(a^{(i)}\right),
    \quad \bar{g} = \frac{1}{n}\sum_{j=1}^{n} g^{(j)} .
  \]
\end{enumerate}

Após $T$ iterações, o vetor de conceito $v$ é tomado como o autovetor principal de $M_T$,
com $\|v\|_2 = 1$. Ele é então orientado calculando-se a correlação de Pearson $\rho$
entre as projeções $\{\langle a^{(i)}, v\rangle\}$ e os rótulos $\{y^{(i)}\}$ e
substituindo-se $v$ por $\operatorname{sign}(\rho)\, v$, de forma que uma projeção positiva
signifique ``mais conceito''. O gradiente que move a AGOP, escrevendo
$d_i(z) = \sqrt{(z - a^{(i)})^{\top} M (z - a^{(i)})}$, é então

\[
  \nabla_z f_t(z) \;=\; -\frac{1}{\sigma}\sum_{i=1}^{n}
    \alpha_{t,i}\, K_{M_t}\!\left(a^{(i)}, z\right)\,
    \frac{1}{d_i(z)}\; M\left(z - a^{(i)}\right),
\]

com a convenção de que os termos para os quais $d_i(z) = 0$ são descartados. Os
hiperparâmetros típicos são $\lambda = 10^{-3}$, $\sigma \in \{1, 10, 100\}$ e
$T \leq 10$.

Preferimos a RFM a uma regressão logística simples porque a AGOP reporta as direções às
quais um preditor \emph{não linear} é sensível, de modo que ela pode, em princípio, trazer
à tona estrutura conceitual que um único separador linear não enxerga. Beaglehole et al.
sustentam essa afirmação sob duas condições, a saber, conceitos abstratos e modelos
grandes, e vale manter as condições atreladas a ela, uma vez que são elas que decidem o
desfecho da §\ref{sec:baselines-anchor} adiante, onde, sobre os três conceitos fáceis e
dentro do domínio disponíveis aqui, a LogReg iguala e supera ligeiramente a
RFM.

\subsection{\textit{Steering}}

Dados os vetores de conceito $\{v_l\}$ obtidos por bloco, aplicamos \textit{steering} ao
modelo durante a passagem direta somando $\varepsilon v_l$ a cada linha de \textit{token} da saída do bloco
$l$, isto é,

\[
  \tilde{A}_{\ell,i}(X) = A_{\ell,i}(X) + \varepsilon\, v_\ell
  \qquad \text{para toda posição de token } i,
\]

onde $\varepsilon$ é o coeficiente de controle e o vetor de conceito é tomado com
$\|v_\ell\|_2 = 1$. Note que o sinal de $\varepsilon$ carrega a direção da intervenção,
uma vez que um $\varepsilon$ positivo empurra o modelo em direção ao conceito e um negativo
o empurra para longe dele. Caso se deseje aplicar \textit{steering} a um grupo de conceitos
simultaneamente, deve-se utilizar uma combinação dos vetores correspondentes.

\subsection{\textit{Monitoring}}

Para o \textit{monitoring}, ao invés do \textit{steering}, guardamos os $p$ autovetores principais
$\{v_{l,j}\}$ por bloco, e não um único deles. A característica do \textit{prompt} $i$ é
então o conjunto das projeções $\langle a_l^{(i)}, v_{l,j} \rangle$ sobre todos os blocos
$l$ e todas as componentes $j$, o que constitui um vetor em $\mathbb{R}^{Lp}$. Sobre essas
características treinamos um classificador sob uma de duas estratégias, ou por bloco,
treinando um classificador para cada bloco e retendo o melhor deles, ou agregada,
concatenando todos os blocos e treinando um único classificador. A avaliação é feita pela
AUROC (\textit{Area Under the Receiver Operating Characteristic curve}, ou área sob a
curva característica de operação do receptor).

\subsection{\textit{Baselines}}

A RFM é comparada contra três métodos de extração mais baratos, descritos a seguir. Os
três operam sobre as ativações de treinamento centradas na média, retornam um vetor
unitário e são orientados pela mesma regra do sinal de Pearson usada na RFM, de forma que
a comparação entre eles seja uma comparação de direções e de mais nada.

\begin{itemize}
  \item \textbf{PCA.} Pareia-se a $j$-ésima amostra positiva com a $j$-ésima amostra
  negativa, $j = 1, \dots, m$ com $m = \min(n_+, n_-)$, e forma-se a matriz de diferenças
  $D$ com linhas $D_j = a^{(\pi^+(j))} - a^{(\pi^-(j))}$. Centra-se essa matriz na média,
  $\tilde{D} = D - \bar{D}$, de modo que $v$ seja o vetor singular à direita principal de
  $\tilde{D}$, equivalentemente o autovetor principal de $\tilde{D}^{\top}\tilde{D}$.

  \item \textbf{DiffMeans} (diferença de médias).
  $v = \dfrac{\mu_1 - \mu_0}{\left\|\mu_1 - \mu_0\right\|_2}$, onde
  $\mu_c = \operatorname{mean}(S_c)$ é a ativação média da classe $c$.

  \item \textbf{LogReg} (regressão logística). $v = w / \|w\|_2$, para o vetor de coeficientes
  ajustado $w$, com a força de regularização inversa $C$ varrida na partição de validação.
\end{itemize}
